\documentclass[12pt, twoside]{article}
\input{../preamble}

\fancyhead[LE]{\thepage}
\fancyhead[RO]{Markscheme}
\fancyhead[LO]{La Scuola d'Italia / Dr.\ Huson / IB Math 11 \\* 16 September 2026}

\begin{document}
\subsubsection*{1.4 Pre-Quiz: Arithmetic Sequences --- Markscheme}

\emph{Calculator allowed. Marks are not printed on the student sheet; they are
recorded here so the set can be scored consistently. Total: 54 marks.}

\begin{enumerate}[itemsep=0.3cm]

%--- Problem 1: Arithmetic vs non-arithmetic; the constant-difference test [9 marks] ---
\item[1.] [Maximum mark: 9]

\medskip

\textbf{(a)} The three differences $u_{2} - u_{1}$, $u_{3} - u_{2}$,
$u_{4} - u_{3}$:

\begin{center}
\begin{tabular}{rll}
(i) & $6,\; 10,\; 14,\; 18,\; 22$ & $4,\; 4,\; 4$ \\
(ii) & $3,\; 6,\; 12,\; 24,\; 48$ & $3,\; 6,\; 12$ \\
(iii) & $1,\; 4,\; 9,\; 16,\; 25$ & $3,\; 5,\; 7$ \\
(iv) & $20,\; 14,\; 8,\; 2,\; -4$ & $-6,\; -6,\; -6$ \\
(v) & $5,\; 5,\; 5,\; 5,\; 5$ & $0,\; 0,\; 0$ \\
\end{tabular}
\end{center}
\hfill \textbf{A1A1}

\emph{A1 for at least three of the five rows correct; A1 for all five. The
negative signs in (iv) and the zeros in (v) must be present.}

\medskip

\textbf{(b)} Arithmetic: (i), (iv) and (v). \hfill \textbf{A1}

$d = 4$ in (i), $d = -6$ in (iv) \hfill \textbf{A1}

$d = 0$ in (v) \hfill \textbf{A1}

\emph{The third A1 is for treating $d = 0$ as a legitimate common difference.
A constant sequence is arithmetic; a candidate who excludes (v) earns the
first two marks only.}

\emph{FT: award the $d$ marks for values that follow correctly from the
candidate's own differences in (a).}

\medskip

\textbf{(c)} (ii) is not arithmetic: each term is the previous term
\emph{multiplied} by $2$. \hfill \textbf{A1}

(iii) is not arithmetic: these are the square numbers, $u_{n} = n^{2}$; the
differences grow by $2$ each time. \hfill \textbf{A1}

\emph{Accept ``it doubles'' or ``geometric with ratio 2'' for (ii). Accept
``squares'', ``$1^{2}, 2^{2}, 3^{2}, \ldots$'', or ``the differences are the
odd numbers'' for (iii).}

\medskip

\textbf{(d)} ``A sequence is arithmetic when the difference between
consecutive terms \ldots \hfill \textbf{A1}

\ldots\ is the same (constant) for every pair of consecutive terms.''
\hfill \textbf{A1}

\emph{A1 for referring to the difference between consecutive (successive,
neighbouring) terms, A1 for saying that difference is constant. ``You add the
same number each time'' earns both. ``It goes up by the same amount'' earns
both. ``There is a common difference'' alone earns the second mark only ---
it names the term without saying what it means.}

\newpage

%--- Problem 2: Sequence vs series; u_n vs S_n notation [9 marks] ---
\item[2.] [Maximum mark: 9]

\medskip

\textbf{(a)} A \emph{sequence} is an ordered list of terms.
\hfill \textbf{A1}

A \emph{series} is what you get when the terms of a sequence are added
together. \hfill \textbf{A1}

\emph{Accept any wording that carries the two ideas: a sequence lists,
a series sums. ``A series has plus signs'' earns the second mark.}

\medskip

\textbf{(b)} $4 + 9 + 14 + 19 + 24$ \hfill \textbf{A1}

$= 70$ \hfill \textbf{A1}

\medskip

\textbf{(c)} $u_{5}$ is the fifth term of the sequence; $u_{5} = 24$.
\hfill \textbf{A1}

$S_{5}$ is the sum of the first five terms; $S_{5} = 70$.
\hfill \textbf{A1}

\emph{Each A1 needs both the meaning and the value. A value with no
description, or a description with no value, earns nothing for that part.}

\emph{FT: award the second A1 for the candidate's own sum from (b).}

\medskip

\textbf{(d)} $3,\; 7,\; 11,\; 15$ --- sequence \qquad
$3 + 7 + 11 + 15$ --- series

$u_{n} = 2n + 1$ --- sequence \qquad $S_{20} = 430$ --- series
\hfill \textbf{A1A1}

\emph{A1 for any two correct, A1 for all four. $u_{n} = 2n + 1$ is a rule for
the terms of a sequence; $S_{20}$ is the value of a series.}

\medskip

\textbf{(e)} $S_{5} = u_{1} + u_{2} + u_{3} + u_{4} + u_{5}$ and
$S_{4} = u_{1} + u_{2} + u_{3} + u_{4}$, so subtracting removes the first
four terms and leaves $u_{5}$. \hfill \textbf{R1}

\emph{Accept the general statement $S_{n} - S_{n-1} = u_{n}$ explained in the
same way, or a numerical demonstration on the sequence in (b):
$70 - 46 = 24 = u_{5}$.}

\newpage

%--- Problem 3: The Gauss story; pairing and the S_n = n/2 (u_1 + u_n) formula [9 marks] ---
\item[3.] [Maximum mark: 9]

\medskip

\textbf{(a)} $1 + 100 = 101$, \quad $2 + 99 = 101$, \quad $3 + 98 = 101$
--- every pair has the same total, $101$. \hfill \textbf{A1}

\emph{The A1 is for the observation, not for the three additions. As one
number goes up by 1 the other goes down by 1, so the total does not change.}

\medskip

\textbf{(b)} The $100$ numbers make $\dfrac{100}{2} = 50$ pairs.
\hfill \textbf{M1}

$50 \times 101 = 5050$ \hfill \textbf{A1}

\emph{M1 for $100 \div 2$, or for a clear statement that the pairs run
$1 + 100$ down to $50 + 51$.}

\medskip

\textbf{(c)} $u_{1} = 1$, \quad $d = 1$, \quad $u_{100} = 100$

$S_{100} = \dfrac{100}{2}\bigl(1 + 100\bigr)$ \hfill \textbf{M1}

$= 50 \times 101 = 5050$ \hfill \textbf{A1}

\emph{M1 for substituting all three values into the formula; A1 for $5050$
together with a statement that it agrees with (b).}

\emph{FT: award the A1 for a value correctly computed from the candidate's own
$u_{1}$, $d$ and $u_{100}$.}

\medskip

\textbf{(d)} Pairing the first term with the last, the second with the
second-to-last, and so on gives a total of $u_{1} + u_{n}$ every time.
\hfill \textbf{R1}

There are $\dfrac{n}{2}$ such pairs, because $n$ terms make $n$ divided by
$2$ pairs --- so $S_{n} = \dfrac{n}{2}(u_{1} + u_{n})$ is the number of pairs
times the pair total. \hfill \textbf{R1}

\emph{The second R1 is specifically for identifying $\frac{n}{2}$ as the
number of pairs. Accept the alternative explanation: write the series
forwards and backwards, add the two copies to get $n$ lots of
$(u_{1} + u_{n})$, then halve --- R1 for the doubling set-up, R1 for the
halving.}

\medskip

\textbf{(e)} $\dfrac{500}{2} = 250$ pairs, each totalling $1 + 500 = 501$
\hfill \textbf{M1}

$250 \times 501 = 125\,250$ \hfill \textbf{A1}

\emph{Accept the formula route $S_{500} = \frac{500}{2}(1 + 500)$ for the M1.}

\newpage

%--- Problem 4: Fibonacci; recursive vs explicit rules [9 marks] ---
\item[4.] [Maximum mark: 9]

\medskip

\textbf{(a)} $21,\; 34,\; 55$ \hfill \textbf{A1}

\emph{All three required for the mark.}

\medskip

\textbf{(b)} Differences: $1 - 1 = 0$, \quad $2 - 1 = 1$, \quad
$3 - 2 = 1$, \quad $5 - 3 = 2$, \quad $8 - 5 = 3$ \hfill \textbf{M1}

The differences are not all the same, so the sequence is not arithmetic.
\hfill \textbf{A1}

\emph{M1 for at least three differences computed; A1 for the conclusion
linked to the differences not being constant. A bare ``it is not
arithmetic'' earns nothing.}

\medskip

\textbf{(c)} A rule is \emph{recursive} when each new term is built from the
term or terms that came before it. \hfill \textbf{R1}

A recursive rule also needs starting value(s) --- here $u_{1} = u_{2} = 1$ ---
and to reach a term you must work through all the terms before it.
\hfill \textbf{R1}

\emph{Accept ``the rule refers back to itself'', ``the answer to one step is
the input to the next step''. The second R1 is for either of the two ideas:
that starting values are required, or that you cannot jump straight to
$u_{n}$.}

\medskip

\textbf{(d)} Recursive: $7 \to 11 \to 15 \to 19 \to 23$, so $u_{5} = 23$
\hfill \textbf{A1}

Explicit: $u_{5} = 4(5) + 3 = 23$ \hfill \textbf{A1}

For $u_{100}$, use the explicit rule: it gives the answer in one step,
$u_{100} = 4(100) + 3 = 403$, whereas the recursive rule needs $99$
additions first. \hfill \textbf{R1}

\emph{The first A1 requires the intermediate terms to be shown --- that is
what makes it the recursive route. R1 for choosing the explicit rule with a
reason about the number of steps; the value $403$ is not required.}

\medskip

\textbf{(e)} $\dfrac{u_{10}}{u_{9}} = \dfrac{55}{34} = 1.6176\ldots = 1.62$
(3 s.f.) \quad and \quad
$\dfrac{u_{12}}{u_{11}} = \dfrac{144}{89} = 1.6180\ldots = 1.62$ (3 s.f.);
the ratios are settling down to about $1.618$. \hfill \textbf{A1}

\emph{The mark needs both values and a comment that they are close to each
other or approaching a fixed number. Naming the golden ratio is not
required, but mention it when returning the work.}

\newpage

%--- Problem 5: Recursive self-improvement (RSI) in the AI news; arithmetic vs non-arithmetic growth [9 marks] ---
\item[5.] [Maximum mark: 9]

\medskip

\textbf{(a)} The improvement is applied to its own output: the system that
round $n$ produces is the system that does the improving in round $n + 1$,
exactly as $u_{n+1}$ is defined from $u_{n}$. \hfill \textbf{R1}

\emph{Accept any answer that carries ``the result of one round becomes the
input to the next''. ``It repeats'' alone is not enough --- repetition is not
the same as feeding the output back in.}

\medskip

\textbf{(b)} Model A: $100,\; 120,\; 140,\; 160$ \quad
Model B: $100,\; 120,\; 144,\; 172.8$ \hfill \textbf{A1}

Model A is arithmetic: its differences are $20, 20, 20$ --- constant. Model
B's differences are $20, 24, 28.8$, which are not constant.
\hfill \textbf{R1}

\emph{A1 for both lists. R1 for naming Model A \emph{and} giving a reason
based on the differences.}

\medskip

\textbf{(c)} Model A: $c_{1} = 100$, \quad $c_{n+1} = c_{n} + 20$
\hfill \textbf{A1}

Model B: $c_{1} = 100$, \quad $c_{n+1} = 1.2\,c_{n}$ \hfill \textbf{A1}

\emph{Accept $c_{n+1} = c_{n} + 0.2c_{n}$ for Model B. The starting value
must appear in at least one of the two rules; do not withhold both marks if
it is stated once for both models.}

\medskip

\textbf{(d)} Model A: $c_{10} = 100 + 9(20) = 280$ \hfill \textbf{A1}

Model B: $c_{10} = 100(1.2)^{9} = 515.97\ldots = 516$ (3 s.f.)
\hfill \textbf{A1}

\emph{Accept $516$, $515.98$, or the full $515.978\ldots$ Accept a correct
value reached by stepping through all ten rounds on the GDC.}

\emph{FT: award each A1 for a value that follows correctly from the
candidate's own rule in (c).}

\medskip

\textbf{(e)} Model B. \hfill \textbf{A1}

Each round adds $20\%$ of a score that is already bigger than last round's,
so the increase itself grows --- $20$, then $24$, then $28.8$ --- and the
model speeds up. Model A adds the same $20$ points for ever.
\hfill \textbf{R1}

\emph{R1 requires the reason to be about the \emph{increase} growing, not
just about Model B being larger. FT: award both marks for a choice argued
correctly from the candidate's own terms in (b).}

\newpage

%--- Problem 6: Arithmetic series calculations [9 marks] ---
\item[6.] [Maximum mark: 9]

$4 + 11 + 18 + \cdots + 200$

\medskip

\textbf{(a)} $d = 7$ \hfill \textbf{A1}

\medskip

\textbf{(b)} $u_{n} = 4 + 7(n - 1) = 200$ \hfill \textbf{M1}

$7n - 3 = 200$, so $7n = 203$ and $n = 29$ \hfill \textbf{A1}

\emph{M1 for setting the $n^{\text{th}}$-term expression equal to $200$,
however it is written. Accept a GDC table or list used to locate $200$ as
the $29^{\text{th}}$ term.}

\medskip

\textbf{(c)} $S_{29} = \dfrac{29}{2}\bigl(4 + 200\bigr)$ \hfill \textbf{M1}

$= 29 \times 102 = 2958$ \hfill \textbf{A1}

\emph{Accept the other form, $S_{29} = \frac{29}{2}(2(4) + 28(7))$.}

\emph{FT: award the A1 for a sum computed correctly from the candidate's own
$n$ in (b).}

\medskip

\textbf{(d)} $S_{30} = \dfrac{30}{2}\bigl(2(6) + 29(-2)\bigr)$
\hfill \textbf{M1}

$= 15(12 - 58) = 15(-46) = -690$ \hfill \textbf{A1}

\emph{M1 for substituting $u_{1} = 6$, $d = -2$, $n = 30$ into either sum
formula. The negative answer is correct --- the terms pass through zero and
the later terms are negative.}

\medskip

\textbf{(e)} $S_{10} = \dfrac{10}{2}\bigl(2(5) + 9d\bigr) = 320$
\hfill \textbf{M1}

$5(10 + 9d) = 320$

$10 + 9d = 64$

$9d = 54$, so $d = 6$ \hfill \textbf{A1}

\emph{M1 for a correct equation in $d$; A1 for $d = 6$. Accept the equation
solved on the GDC.}

\end{enumerate}
\end{document}
